\documentclass[11pt]{article}
% ===== Shared preamble for Calendar notes =====
\usepackage[margin=1in]{geometry}
\usepackage{amsmath,amssymb}
\usepackage{enumitem}
\usepackage{xcolor}
\usepackage{tcolorbox}
\tcbuselibrary{skins,breakable}
\usepackage{titlesec}
\usepackage{fancyhdr}
\usepackage{hyperref}
\usepackage{array}
\usepackage{booktabs}
\usepackage{longtable}
\usepackage{pifont} % for \ding

% ---- Colors ----
\definecolor{primary}{HTML}{1B3A6B}
\definecolor{accent}{HTML}{2E7D32}
\definecolor{warnclr}{HTML}{B45309}
\definecolor{tipclr}{HTML}{0F5D8C}
\definecolor{lightbg}{HTML}{F4F6F8}

% ---- Section styling ----
\titleformat{\section}
  {\normalfont\Large\bfseries\color{primary}}
  {\thesection}{0.5em}{}[{\color{primary}\titlerule[1pt]}]
\titleformat{\subsection}
  {\normalfont\large\bfseries\color{primary}}
  {\thesubsection}{0.5em}{}

% ---- Header/Footer ----
\pagestyle{fancy}
\fancyhf{}
\renewcommand{\headrulewidth}{0.4pt}
\fancyfoot[C]{\thepage}

% ---- Hyperref ----
\hypersetup{
  colorlinks=true,
  linkcolor=primary,
  urlcolor=tipclr,
  citecolor=accent
}

% ---- Custom boxes ----
% Formula / Key concept box
\newtcolorbox{formulabox}[1][]{
  enhanced, breakable,
  colback=lightbg, colframe=primary,
  boxrule=0.8pt, arc=2pt,
  left=8pt, right=8pt, top=6pt, bottom=6pt,
  title={\textbf{#1}},
  fonttitle=\bfseries\color{white},
  colbacktitle=primary,
  attach boxed title to top left={xshift=8pt, yshift=-2mm},
  boxed title style={boxrule=0pt, arc=2pt}
}

% Tip / trick box
\newtcolorbox{tipbox}[1][Tip]{
  enhanced, breakable,
  colback=blue!4, colframe=tipclr,
  boxrule=0.8pt, arc=2pt,
  left=8pt, right=8pt, top=6pt, bottom=6pt,
  title={\textbf{\textcolor{white}{\ding{72}\ #1}}},
  fonttitle=\bfseries, colbacktitle=tipclr,
}

% Warning / common mistake box
\newtcolorbox{warnbox}[1][Watch Out]{
  enhanced, breakable,
  colback=orange!6, colframe=warnclr,
  boxrule=0.8pt, arc=2pt,
  left=8pt, right=8pt, top=6pt, bottom=6pt,
  title={\textbf{\textcolor{white}{#1}}},
  fonttitle=\bfseries, colbacktitle=warnclr,
}

% Example / worked question box
\newtcolorbox{examplebox}[1][Example]{
  enhanced, breakable,
  colback=green!4, colframe=accent,
  boxrule=0.8pt, arc=2pt,
  left=8pt, right=8pt, top=6pt, bottom=6pt,
  title={\textbf{\textcolor{white}{#1}}},
  fonttitle=\bfseries, colbacktitle=accent,
}


\title{\color{primary}\textbf{Calendar --- Fundamentals}\\[4pt]\large Odd Days, Leap Years \& the Day-of-the-Week Formula}
\author{Aptitude Grind}
\date{}

\lhead{\small Calendar --- Fundamentals}
\rhead{\small Aptitude Grind}

\begin{document}
\maketitle
\thispagestyle{fancy}

\section*{Why Calendar Problems Matter}
Calendar questions test whether you can find the \textbf{day of the week} for a given date, or work
backward/forward from a known day. They show up constantly in placement tests (TCS NQT, Infosys,
Wipro, Capgemini) and in CAT/other reasoning sections, usually as 1--2 quick-scoring questions
\textbf{if} you know the odd-days method cold. The entire topic rests on one idea: \emph{the calendar
repeats every 7 days, so only the remainder after dividing by 7 matters.}

\section{Odd Days}

\begin{formulabox}[Definition]
In a given period, the number of days \textbf{more than complete weeks} is called the number of
\textbf{odd days}.
\end{formulabox}

For example, 17 days = 2 complete weeks (14 days) + 3 extra days $\Rightarrow$ 3 odd days.
Only the odd-day count matters for finding the day of the week --- full weeks never shift the day.

\subsection{Odd Days in a Year}
\begin{itemize}[leftmargin=1.4em]
  \item 1 ordinary year $= 365$ days $= 52$ weeks $+ 1$ day $\Rightarrow$ \textbf{1 odd day}
  \item 1 leap year $= 366$ days $= 52$ weeks $+ 2$ days $\Rightarrow$ \textbf{2 odd days}
\end{itemize}

\subsection{Odd Days in Centuries}
\begin{formulabox}[Century Odd-Day Table --- memorize this]
\centering
\begin{tabular}{lc}
\toprule
\textbf{Period} & \textbf{Odd Days} \\
\midrule
100 years & 5 \\
200 years & 3 \\
300 years & 1 \\
400 years & 0 \\
\bottomrule
\end{tabular}
\end{formulabox}

\textbf{Why 400 years gives 0:} $100$ years $= 76$ ordinary $+ 24$ leap years
$= (76 \times 1 + 24 \times 2) = 124$ odd days $= 17$ weeks $+ 5$ days $\Rightarrow 5$ odd days.
So $400$ years $= (5 \times 4) + 1 \text{ extra leap day (400 itself is a leap century)} = 21 \equiv 0 \pmod 7$.

\begin{tipbox}[Pattern]
Every 4th century (400, 800, 1200, 1600, 2000, 2400, \ldots) has \textbf{0 odd days} --- meaning the
calendar of year $Y$ is identical to the calendar of year $Y + 400$. This is why 2000 and 2400 have
the exact same calendar.
\end{tipbox}

\section{Leap Years}

\begin{formulabox}[Leap Year Rule]
\begin{enumerate}[leftmargin=1.4em]
  \item Every year divisible by 4 is a leap year --- \textbf{unless} it is a century year (ends in 00).
  \item A century year is a leap year only if it is divisible by \textbf{400}.
\end{enumerate}
\end{formulabox}

\begin{examplebox}[Quick Classification]
\begin{itemize}[leftmargin=1.4em]
  \item \textbf{Leap:} 1948, 2004, 1676, 2000, 1600, 2400 (divisible by 4; century years also divisible by 400)
  \item \textbf{Not leap:} 2001, 2002, 2003, 2005 (not divisible by 4); 1800, 1900, 2100, 2300
    (century years \emph{not} divisible by 400)
\end{itemize}
\end{examplebox}

\begin{warnbox}[Common Mistake]
Students often forget the century exception and assume any year divisible by 4 is automatically a
leap year. \textbf{1900 is divisible by 4 but is NOT a leap year} (not divisible by 400). Always check
century years separately.
\end{warnbox}

\section{Mapping Odd Days to a Day of the Week}

Once you compute the total number of odd days for a period (taking the result mod 7), map it using:

\begin{formulabox}[Odd Days $\to$ Day Table]
\centering
\begin{tabular}{lccccccc}
\toprule
Odd Days & 0 & 1 & 2 & 3 & 4 & 5 & 6 \\
\midrule
Day & Sun & Mon & Tue & Wed & Thu & Fri & Sat \\
\bottomrule
\end{tabular}
\end{formulabox}

This table is anchored on the fact that \textbf{1 January, 1 AD was a Monday} in the proleptic
Gregorian calendar used for these problems, giving Sunday = 0 odd days as the baseline.

\section{General Method: Finding the Day for Any Date}

To find the day of the week on date $D$ of month $M$, year $Y$:

\begin{formulabox}[Step-by-Step Method]
\begin{enumerate}[leftmargin=1.4em]
  \item Split $Y - 1$ (the reference is "up to the end of the previous year") into:
    \begin{itemize}
      \item Number of 400s $\to$ always 0 odd days
      \item Remaining 100s (0 to 3) $\to$ 5, 3, or 1 odd days respectively
      \item Remaining years $\to$ (ordinary years $\times 1$) + (leap years $\times 2$)
    \end{itemize}
  \item Add the odd days contributed by the \textbf{months elapsed so far in year $Y$} (Jan 1 to the
    last day of the previous month), remembering to use 29 for February in a leap year.
  \item Add the day number $D$ of the current month.
  \item Sum everything, take mod 7, and read off the day from the Odd Days $\to$ Day table.
\end{enumerate}
\end{formulabox}

\textbf{Don't split the years into 1600 + 400 + leftover by hand.} That decomposition appears in a
lot of older textbooks, but it's a slow way of doing one thing: throwing away every complete
400-year block, because a complete 400-year block always contributes exactly 0 odd days. There is a
one-line way to get the same ``leftover'' number:

\begin{tipbox}[Shortcut]
Odd days from complete years up to year $Y-1$ depend only on $(Y-1) \bmod 400$, because every full
400-year block cancels out. Take the previous year, reduce it mod 400, and count odd days for
\emph{only that many years} --- no need to think about which specific 400-year block you're in.
\end{tipbox}

\begin{examplebox}[Worked Example --- 15 August 2010]
\textbf{Step 1: Years up to 2009.}
\[
2009 \bmod 400 = 9 \quad (2000 = 5\times 400,\ \text{remainder } 9)
\]
So you only need odd days for the 9 years \textbf{2001--2009} --- nothing before that matters.
\begin{itemize}[leftmargin=1.4em]
  \item Leap years in 2001--2009: 2004, 2008 $\Rightarrow$ 2 leap, 7 ordinary
  \item Odd days $= (2 \times 2) + (7 \times 1) = 4 + 7 = 11 \equiv \textbf{4} \pmod 7$
\end{itemize}
\textbf{Step 2: Months elapsed (Jan--Jul 2010) + 15 days of Aug.}
\[
31+28+31+30+31+30+31+15 = 227 \text{ days} = 32 \text{ weeks} + 3 \text{ days} \Rightarrow 3 \text{ odd days}
\]
\textbf{Step 3: Total.}
\[
4 + 3 = 7 \equiv 0 \pmod 7 \Rightarrow \textbf{Sunday}
\]
15 August 2010 was indeed a \textbf{Sunday}.
\end{examplebox}

\begin{warnbox}[Common Mistake]
Writing $2009 = 1600 + 400 + 9$ and separately zeroing out the 1600 and the 400 blocks. It reaches
the same right answer, but it's two unnecessary steps under exam time pressure. $2009 \bmod 400 = 9$
gets you the same ``9 leftover years'' in one line --- that's all the extra decomposition was ever
doing.
\end{warnbox}

\begin{tipbox}[Faster Alternative]
For placement-test speed, most people skip the ``years since year 1'' method entirely and instead
work \textbf{relative to a known anchor date} (e.g., ``1 Jan 2000 was a Saturday'') using only the
years \emph{between} the anchor and the target date --- no mod-400 needed at all, since the anchor is
already recent. See \texttt{02-fast-tricks-and-shortcuts.tex} for the anchor-date and month-code
methods, and \texttt{03-question-pattern-catalog.tex} for every question type you'll meet in
placement tests, each with its own fastest method.
\end{tipbox}

\section{Key Takeaways}
\begin{itemize}[leftmargin=1.4em]
  \item Everything reduces to counting \textbf{odd days} and taking the result mod 7.
  \item Memorize: 100y $\to$ 5, 200y $\to$ 3, 300y $\to$ 1, 400y $\to$ 0 odd days.
  \item Leap year = divisible by 4, except century years, which need divisibility by 400.
  \item Always double check whether February in the target year has 28 or 29 days.
\end{itemize}

\end{document}
